\documentclass[sigconf]{acmart}

%
% \BibTeX command to typeset BibTeX logo in the docs
\AtBeginDocument{%
	\providecommand\BibTeX{{%
			Bib\TeX}}}



\copyrightyear{2024}
\acmYear{2024}
\setcopyright{rightsretained}
\acmConference[CCS '24]{Proceedings of the 2024 ACM SIGSAC Conference on Computer and Communications Security}{October 14--18, 2024}{Salt Lake City, UT, USA}
\acmBooktitle{Proceedings of the 2024 ACM SIGSAC Conference on Computer and Communications Security (CCS '24), October 14--18, 2024, Salt Lake City, UT, USA}
\acmDOI{10.1145/3658644.3670356}
\acmISBN{979-8-4007-0636-3/24/10}

\makeatletter
\gdef\@copyrightpermission{
	\begin{minipage}{0.3\columnwidth}
		\href{https://creativecommons.org/licenses/by/4.0/}{\includegraphics[width=0.90\textwidth]{4ACM-CC-by-88x31.eps}}
	\end{minipage}\hfill
	\begin{minipage}{0.7\columnwidth}
		\href{https://creativecommons.org/licenses/by/4.0/}{This work is licensed under a Creative Commons Attribution International 4.0 License.}
	\end{minipage}
	\vspace{5pt}
}
\makeatother



\usepackage{amsmath, amsthm,amsfonts,enumitem,tikz, subcaption, mdframed}
\usepackage{multirow}
%\usepackage{amssymb}
\usepackage{tcolorbox}
\usepackage{comment}
\usepackage{graphicx}
\usepackage{caption}
\usepackage{subcaption}
\usepackage{hyperref}
\usepackage{pgfplots}
\usepackage{booktabs, tabularx}
\usepackage{colortbl}
\usepackage{url}
\usetikzlibrary{arrows,arrows.meta,shapes.arrows, calc, positioning,matrix}



\begin{document}
	
	
	\title{Batching-Efficient RAM using Updatable Lookup Arguments}
	
	
	
	\author{Moumita Dutta}
	\affiliation{
		\institution{Indian Institute of Science}
		\city{Bangalore}
		\country{India}}
	\email{moumitadutta@iisc.ac.in}
	
	\author{Chaya Ganesh}
	\affiliation{
		\institution{Indian Institute of Science}
		\city{Bangalore}
		\country{India}}
	\email{chaya@iisc.ac.in}
	
	\author{Sikhar Patranabis}
	\affiliation{
		\institution{IBM Research India}
		\city{Bangalore}
		\country{India}}
	\email{sikhar.patranabis@ibm.com}
	
	\author{Shubh Prakash}
	\affiliation{
		\institution{Indian Institute of Science}
		\city{Bangalore}
		\country{India}}
	\email{shubhprakash@iisc.ac.in}
	
	\author{Nitin Singh}
	\affiliation{
		\institution{IBM Research India}
		\city{Bangalore}
		\country{India}}
	\email{nitisin1@in.ibm.com}
	
	
	
	
	
	\begin{abstract}		
		This document contains an artifact appendix that summarizes the purpose, availability, content, dependencies, and execution of the artifacts that support our paper. The camera-ready version of our paper is attached separately as part of the artifact submission. 
		
	\end{abstract}
	
	\maketitle
	
	
	
	\begin{CCSXML}
		<ccs2012>
		<concept>
		<concept_id>10002978.10002979.10002981</concept_id>
		<concept_desc>Security and privacy~Public key (asymmetric) techniques</concept_desc>
		<concept_significance>500</concept_significance>
		</concept>
		</ccs2012>
	\end{CCSXML}
	%	
	%	
	\ccsdesc[500]{Security and privacy~Public key (asymmetric) techniques}
	
	\keywords{Succinct Arguments, Efficient RAM, Indexed lookups, Lookup Arguments, Rollups}
	
	
	
	%%%%%%%%%%%%%%%%%%%%%%%%%%%%%%%%%%%%%%%%%%%%%%%%%%%%%%%%%%%%%%%%%%%%%
	% Artifact Appendix Template for ACM CCS 2024.
	% See end of this file for authorship and provenance info.
	%
	% Please limit your Artifact Appendix to no more than two pages.
	%%%%%%%%%%%%%%%%%%%%%%%%%%%%%%%%%%%%%%%%%%%%%%%%%%%%%%%%%%%%%%%%%%%%%
	
	%% This template uses list-making environments from the `paralist`
	%% LaTeX package.  Add the following to the preamble of your LaTeX
	%% document (somewhere before `\begin{document}`:
	%%
	%% \usepackage{paralist}
	
	%%%%%%%%%%%%%%%%%%%%%%%%%%%%%%%%%%%%%%%%%%%%%%%%%%%%%%%%%%%%%%%%%%%%%
	
	
	
	\appendix
	
	\section{Artifact Appendix}
	
	\emph{The Artifact Appendix is meant to summarize the purpose,
		availability, content, dependencies, and execution of the artifacts
		that support your paper.  It should include a clear description of
		your artifacts' hardware, software, and configuration requirements.
		If your artifacts have received the \emph{Artifacts
			Evaluated--Functional}, \emph{Artifacts Evaluated--Reusable}, or
		\emph{Results Reproduced} badge, this appendix should summarize the
		relevant major claims made by your paper and provide instructions for
		validating each claim through the use of your artifacts.  Linking the
		claims of your paper to the artifacts allows a reader of your paper
		to more easily try to reproduce your paper's results.}
	
	\emph{Please fill in all the mandatory sections, keeping their titles
		and organization but removing the current illustrative content, and
		remove the optional sections that do not apply to your artifacts.}
	
	%%%%%%%%%%%%%%%%%%%%%%%%%%%%%%%%%%%%%%%%%%%%%%%%%%%%%%%%%%%%%%%%%%%%%
	
	\subsection{Abstract}
	
	\emph{[Mandatory]}
	%
	\emph{Provide a short description of your artifacts.  What do they
		contain?  How are they packaged?  What do they do when executed?  What
		results to they produce?  How are your artifacts licensed?  Be brief;
		subsequent sections of this appendix can be used fill in the details.}
	
	%%%%%%%%%%%%%%%%%%%%%%%%%%%%%%%%%%%%%%%%%%%%%%%%%%%%%%%%%%%%%%%%%%%%%
	
	\subsection{Description \& Requirements}
	
	\emph{This section should further detail your artifacts and summarize
		the requirements that must be met for a person to recreate an
		experimental environment for utilizing your artifacts.  Where
		applicable, state the minimum hardware and software requirements for
		running your artifacts.  It is also good practice to list and describe
		in this section benchmarks where those are part of, or simply have
		been used to produce results with, your artifacts.}
	
	\emph{[Mandatory]}
	%
	\emph{Here, provide additional details about the general organization
		of your artifacts.  For example, provide a general roadmap to the
		files within your artifacts.  Where is the \texttt{README}?  Where are
		the source files, the datasets, the scripts for running experiments,
		etc.?  You don't need to describe every file and directory here; just
		provide a general ``lay of the land.''}
	
	%%%%%
	
	\subsubsection{Security, privacy, and ethical concerns}
	
	\emph{[Mandatory]}
	%
	\emph{Describe any risks that people might be exposed to while
		executing your artifacts: e.g., risks to their machines' security,
		data privacy, or other ethical concerns.  This is particularly
		important if destructive steps are taken or security mechanisms are
		disabled during execution.}
	
	%%%%%
	
	\subsubsection{How to access}
	
	\emph{[Mandatory]}
	%
	\emph{Describe how to access your artifacts.  If your artifacts have
		received the \emph{Artifacts Available} badge, you must provide the
		DOI(s) for the permanently and publicly available copies of your
		artifacts.  Most likely, this is the version of your artifacts that
		you deposited with Zenodo.}
	
	\emph{You may describe more than one way to access your artifacts.
		For example, if the archived versions of your artifacts are available
		at Zenodo, and you are also making maintained versions of your
		artifacts available through GitHub, you can describe how to access
		both versions of your artifacts.}
	
	%%%%%
	
	\subsubsection{Hardware dependencies}
	
	\emph{[Mandatory]}
	%
	\emph{State any specific hardware features that are required to make
		use of your artifacts: e.g., vendor, CPU/GPU/FPGA, number of
		processors/cores, microarchitecture, interconnect, memory, hardware
		counters, etc.  If your artifacts do not have significant hardware
		dependencies, simply write ``None'' in this section.}
	
	%%%%%
	
	\subsubsection{Software dependencies}
	
	\emph{[Mandatory]}
	%
	\emph{State any specific OS and software packages that are required to
		make use of your artifacts.  This is particularly important if you
		share your source code and it must be compiled, or if your artifacts
		rely on proprietary software that is not included in your artifact
		packages.  If your artifacts do not have significant software
		dependencies, simply write ``None'' in this section.}
	
	%%%%%
	
	\subsubsection{Benchmarks}
	
	\emph{[Mandatory]}
	%
	\emph{Describe any data (e.g., datasets, models, workloads,
		etc.)\ required by the experiments that are reported in your paper and
		supported by your artifacts.  If this does not apply to your
		artifacts, simply write ``None'' in this section.}
	
	%%%%%%%%%%%%%%%%%%%%%%%%%%%%%%%%%%%%%%%%%%%%%%%%%%%%%%%%%%%%%%%%%%%%%
	
	\subsection{Set Up}
	
	\emph{This section should summarize the installation and configuration
		steps that a person should follow to prepare an environment for making
		use of your artifacts.  If the set-up steps are complicated, point the
		reader to the places in your artifacts where detailed set-up
		instructions can be found.}
	
	%%%%%
	
	\subsubsection{Installation}
	
	\emph{[Mandatory]}
	%
	\emph{Provide instructions for downloading and installing dependencies
		as well as the main artifacts.  After following these steps, a user of
		your artifacts should be able to run a simple functionality test.}
	
	%%%%%
	
	\subsubsection{Basic test}
	
	\emph{[Mandatory]}
	%
	\emph{Provide instructions for running a simple functionality test.
		This test does not need to exercise all the features of your
		artifacts, but ideally, it should allow a user to check that all of
		the required software components are correctly installed and
		functioning as intended.  Please include the expected successful
		output and any required input parameters.}
	
	%%%%%%%%%%%%%%%%%%%%%%%%%%%%%%%%%%%%%%%%%%%%%%%%%%%%%%%%%%%%%%%%%%%%%
	
	\subsection{Evaluation Workflow}
	
	\emph{This section should include all the operational steps and
		experiments that a person must be carry out to utilize in your
		artifacts toward the goal of validating your paper's key results and
		claims.  To that end, we ask you to use the two following subsections
		and cross-reference the items therein as explained below.}
	
	%%%%%
	
	\subsubsection{Major claims}
	
	\emph{[Mandatory for \emph{Artifacts Evaluated--Functional},
		\emph{Artifacts Evaluated--Reusable}, and \emph{Results Reproduced};
		optional for \emph{Artifacts Available}]}
	%
	\emph{Enumerate here the major claims (Cx) made in your paper.  Follow
		the examples below.}
	\bigskip
	
	\begin{itemize}
		\item[(C1):]
		%
		\emph{System\_name achieves the same accuracy as the
			state-of-the-art systems for a task X while requiring Y\% less
			storage.  This is proven by the experiment (E1) described in [refer
			to your paper's sections], whose results are illustrated/reported
			in [refer to your paper's plots, tables, sections, etc.].}
		
		\item[(C2):]
		%
		\emph{System\_name has been used to uncover new bugs in the X
			software.  This is proven by experiments (E2) and (E3) in [refer to
			your paper's sections].}
	\end{itemize}
	
	%%%%%
	
	\subsubsection{Experiments}
	
	\emph{[Mandatory for \emph{Artifacts Evaluated--Functional},
		\emph{Artifacts Evaluated--Reusable}, and \emph{Results Reproduced};
		optional for \emph{Artifacts Available}]}
	%
	\emph{Explicitly link the descriptions of your experiments to the
		items you have provided in the previous subsection about major claims.
		Please provide estimates of the person- and compute-time required for
		each of the listed experiments (using the suggested hardware/software
		configuration).  Follow the examples below.}
	
	\bigskip
	\begin{itemize}
		\item[(E1):]
		%
		\emph{[Optional Name] [30~person-minutes + 1~compute-hour + 5\,GB
			disk]: Provide a short explanation of the experiment and expected
			results.}
		
		\emph{Provide the steps to perform the experiment and collect and
			organize the results as expected from your paper.  We encourage you
			to use the following structure with three main blocks for the
			description of your experiment.  If the procedures are complicated,
			point the reader to the places in your artifacts where the full
			instructions can be found.}
		
		\begin{description}
			\item[Preparation:] \emph{Describe the steps required to prepare and
				configure the environment for this experiment.}
			
			\item[Execution:] \emph{Describe the steps required to run this
				experiment.}
			
			\item[Results:] \emph{Describe the steps required to collect and
				interpret the results for this experiment.}
		\end{description}
		
		\item[(E2):] \emph{[Optional Name] [1~person-hour + 3~compute-hours]: \ldots}
		
		\item[(E3):] \emph{[Optional Name] [1~person-hour + 3~compute-hours]: \ldots}
	\end{itemize}
	\bigskip
	
	\emph{In all of the above blocks, please provide indications about the
		expected outcome for each of the steps (given the suggested hardware
		and software configuration).}
	
	%%%%%%%%%%%%%%%%%%%%%%%%%%%%%%%%%%%%%%%%%%%%%%%%%%%%%%%%%%%%%%%%%%%%%
	
	\subsection{Notes on Reusability}
	
	\emph{[Optional]}
	%
	\emph{This section is meant to provide additional information about
		how a person could utilize your artifacts beyond the research
		presented in your paper.  A broad objective of artifact evaluation is
		to encourage you to make your research artifacts reusable by others.}
	
	\emph{Include in this section any instructions, guidance, or advice
		that you believe would help others reuse your artifacts.  For example,
		you could describe how one might scale certain components of your
		artifacts up or down; apply your artifacts to different kinds of
		inputs or datasets; customize your artifacts' behavior by replacing
		specific modules or algorithms; etc.}
	
	%%%%%%%%%%%%%%%%%%%%%%%%%%%%%%%%%%%%%%%%%%%%%%%%%%%%%%%%%%%%%%%%%%%%%
	
	\subsection{Version}
	%%%%%%%%%%%%%%%%%%%%
	% Obligatory.
	% Do not change/remove.
	%%%%%%%%%%%%%%%%%%%%
	Based on the LaTeX template for Artifact Evaluation V20220926.
	
	%%%%%%%%%%%%%%%%%%%%%%%%%%%%%%%%%%%%%%%%%%%%%%%%%%%%%%%%%%%%%%%%%%%%%
	
	\bibliographystyle{plain}
	\bibliography{bib/abbrev3,bib/crypto,bib/other}
	
\end{document}